\subsection{某些表示的连续性}

命题 1. — 设 H 是局部紧群。设 \(\varrho_{1}\)（分别为 \(\varrho_{2}\)）是 G（分别为 H）在分离局部凸 K-向量空间 E \(_{1}\)（分别为 E \(_{2}\)）中的连续表示。令 F 表示赋予紧收敛拓扑的空间 \(\mathcal{L}(\mathrm{E}_{1};\mathrm{E}_{2})\)。由 \(\varrho\) 通过 \(\varrho(g,h)u=\varrho_{2}(h)\circ u\circ\varrho_{1}(g^{-1})\)（\((g,h)\in\mathrm{G}\times\mathrm{H}\)）定义的 G×H 到 F 的表示是连续的。

记 \(\mathcal{L}_c(\mathrm{E}_1;\mathrm{E}_2)\) 为赋予紧收敛拓扑的空间 \(\mathcal{L}(\mathrm{E}_1;\mathrm{E}_2)\)。设 \(\mathfrak{F}_1\)（分别为 \(\mathfrak{F}_2,\mathfrak{F}_3\)）是趋于 \(e\) 的 G 中滤子（分别为趋于 0 的 \(\mathcal{L}_c(\mathrm{E}_1;\mathrm{E}_2)\) 中滤子、趋于 \(e\) 的 H 中滤子）。由于 G 和 H 局部紧，存在 \(\mathrm{C}\in \mathfrak{F}_1\) 和 \(\mathrm{D}\in \mathfrak{F}_3\)，它们是相对紧的。集合 \(\varrho_1(\mathrm{C}^{-1})\) 在 \(\mathcal{L}(\mathrm{E}_1)\) 中等度连续（参见 INT, VIII, p. 129, § 2, no. 1, 注 2, \(a'\)）。由 EVT, III, p. 33 的命题 9 和 EVT, III, p. 31 的命题 4，由 \((u,v)\mapsto u\circ v\) 定义的从 \(\mathcal{L}_c(\mathrm{E}_1;\mathrm{E}_2)\times \varrho_1(\mathrm{C}^{-1})\) 到 \(\mathcal{L}_c(\mathrm{E}_1;\mathrm{E}_2)\) 的映射是连续的。因此，滤基 \(\mathfrak{F}_2\circ \varrho_1(\mathfrak{F}_1^{-1})\) 在 \(\mathcal{L}_c(\mathrm{E}_1;\mathrm{E}_2)\) 中趋于 0。集合 \(\varrho_2(\mathrm{D})\) 在 \(\mathcal{L}(\mathrm{E}_2)\) 中等度连续（参见 INT, VIII, p. 129, § 2, n° 1, rem. 2），并且对每个 \(x\in \mathrm{E}_2\)，集合 \(\varrho_2(\mathrm{D})x\subset \mathrm{E}_2\) 是相对紧的。因此，赋予紧收敛拓扑的 \(\varrho_2(\mathrm{D})\) 在 \(\mathcal{L}(\mathrm{E}_2)\) 中相对紧（TG, X, p. 18, cor. 1）。由 \((u,v)\mapsto u\circ v\) 定义的从 \(\overline{\varrho_2(\mathrm{D})}\times \mathcal{L}_c(\mathrm{E}_1;\mathrm{E}_2)\) 到 \(\mathcal{L}_c(\mathrm{E}_1;\mathrm{E}_2)\) 的映射，由 EVT, III, p. 33 的命题 9 和 EVT, III, p. 31 的命题 4，是连续的。因此，滤基 \(\varrho (\mathfrak{F}_1\times \mathfrak{F}_3)(\mathfrak{F}_2)\) 在 \(\mathcal{L}_c(\mathrm{E}_1;\mathrm{E}_2)\) 中趋于 0。这就证明了断言。

推论. --- 设 \(\varrho\) 是 G 在分离局部凸 K-向量空间 E 上的连续表示。当 E' 赋予紧收敛拓扑时，逆步表示 \(\breve{\varrho}\) 是连续的。

取 \(H = \{e\}\)、\(\varrho_{1} = \varrho\)，并令 \(\varrho_{2}\) 为 K 上的平凡表示，即由命题得到。

注. --- 当 E' 赋予强拓扑时，逆步表示不一定连续（参见 INT, VIII, p. 191, § 2, 习题 3, d)）。可以证明，如果空间 E 是半自反的，则它是连续的（同上，c)）。

